\subsection*{Divergence}

\begin{tcolorbox}[colback=gray!8, colframe=gray!40, title=\textbf{Divergence Toolkit}]
\textbf{Concept family:} Failure modes for convergence of real sequences.

\medskip
\textbf{Atomic items in this section:}
\begin{itemize}
  \item Definition: Divergent Sequence
  \item Definition: Divergence to Positive Infinity
  \item Definition: Divergence to Negative Infinity
  \item Definition: Oscillatory Sequence
  \item Theorem: Divergence to Infinity Implies Real Divergence
  \item Theorem: Two Subsequential Limits Force Divergence
  \item Theorem: Unbounded Above Sequence Has a Positive-Infinity Subsequence
  \item Theorem: Unbounded Below Sequence Has a Negative-Infinity Subsequence
  \item Theorem: Bounded Divergence Produces Two Subsequential Limits
\end{itemize}
\end{tcolorbox}

\begin{exposition}
Divergence is not a single behavior. A sequence can fail to converge because
it escapes upward, escapes downward, oscillates between incompatible limiting
patterns, or remains bounded while refusing to settle. The purpose of this
section is to separate these failure modes and connect them to subsequences.
\end{exposition}

% ---------------------------------------------------------
% Definition --- Divergent Sequence
% ---------------------------------------------------------

\begin{definitionbox}{Definition (Divergent Sequence)}
\begin{definition}[Divergent Sequence]
\label{def:divergent-sequence}
Let $(x_n)$ be a real sequence. The sequence $(x_n)$ is
\textbf{divergent} if there is no $L\in\mathbb{R}$ such that
$x_n\to L$.
\end{definition}
\end{definitionbox}

\begin{remark*}[Standard quantified statement]
\[
\begin{aligned}
&(x_n)\text{ is divergent}\\
&\quad\Longleftrightarrow\quad
\forall L\in\mathbb{R}\;\exists \varepsilon>0\;\forall N\in\mathbb{N}
\;\exists n\ge N\;\bigl(|x_n-L|\ge\varepsilon\bigr).
\end{aligned}
\]
\end{remark*}

\begin{remark*}[Definition predicate reading]
\[
\operatorname{DivergentSequence}(x_n)
\coloneqq
\forall L\in\mathbb{R}\;\exists \varepsilon>0\;\forall N\in\mathbb{N}
\;\exists n\ge N\;\bigl(|x_n-L|\ge\varepsilon\bigr).
\]
\end{remark*}

\begin{remark*}[Negated quantified statement]
\[
(x_n)\text{ is not divergent}
\Longleftrightarrow
\exists L\in\mathbb{R}\;\forall \varepsilon>0\;\exists N\in\mathbb{N}
\;\forall n\ge N\;\bigl(|x_n-L|<\varepsilon\bigr).
\]
\end{remark*}

\begin{remark*}[Negation predicate reading]
\[
\neg\operatorname{DivergentSequence}(x_n)
\Longleftrightarrow
\exists L\in\mathbb{R}\;\operatorname{ConvergentSequence}(x_n,L).
\]
\end{remark*}

\begin{remark*}[Interpretation]
Divergence is the formal negation of real convergence. It says that every
candidate limit fails at some fixed tolerance on every tail.
\end{remark*}

\begin{remark*}[Dependencies]\hfill
\begin{itemize}
  \item \hyperref[def:sequence]{Sequence}
  \item \hyperref[def:convergent-sequence]{Convergence of a Sequence}
\end{itemize}
\end{remark*}

% ---------------------------------------------------------
% Definition --- Divergence to Positive Infinity
% ---------------------------------------------------------

\begin{definitionbox}{Definition (Divergence to Positive Infinity)}
\begin{definition}[Divergence to Positive Infinity]
\label{def:diverges-to-positive-infinity}
Let $(x_n)$ be a real sequence. The sequence $(x_n)$ \textbf{diverges to}
$+\infty$ if
\[
  \forall M\in\mathbb{R}\;\exists N\in\mathbb{N}\;\forall n\ge N
  \;\bigl(M<x_n\bigr).
\]
\end{definition}
\end{definitionbox}

\begin{remark*}[Standard quantified statement]
\[
x_n\to+\infty
\Longleftrightarrow
\forall M\in\mathbb{R}\;\exists N\in\mathbb{N}\;\forall n\ge N
\;\bigl(M<x_n\bigr).
\]
\end{remark*}

\begin{remark*}[Definition predicate reading]
\[
\operatorname{DivergesToPositiveInfinity}(x_n)
\coloneqq
\forall M\in\mathbb{R}\;\exists N\in\mathbb{N}\;\forall n\ge N
\;\bigl(M<x_n\bigr).
\]
\end{remark*}

\begin{remark*}[Negated quantified statement]
\[
x_n\not\to+\infty
\Longleftrightarrow
\exists M\in\mathbb{R}\;\forall N\in\mathbb{N}\;\exists n\ge N
\;\bigl(x_n\le M\bigr).
\]
\end{remark*}

\begin{remark*}[Interpretation]
Divergence to $+\infty$ means every horizontal threshold is eventually left
below the sequence. It is not convergence to a real number; it is escape from
the real line in the positive direction.
\end{remark*}

\begin{remark*}[Dependencies]\hfill
\begin{itemize}
  \item \hyperref[def:sequence]{Sequence}
\end{itemize}
\end{remark*}

% ---------------------------------------------------------
% Definition --- Divergence to Negative Infinity
% ---------------------------------------------------------

\begin{definitionbox}{Definition (Divergence to Negative Infinity)}
\begin{definition}[Divergence to Negative Infinity]
\label{def:diverges-to-negative-infinity}
Let $(x_n)$ be a real sequence. The sequence $(x_n)$ \textbf{diverges to}
$-\infty$ if
\[
  \forall M\in\mathbb{R}\;\exists N\in\mathbb{N}\;\forall n\ge N
  \;\bigl(x_n<M\bigr).
\]
\end{definition}
\end{definitionbox}

\begin{remark*}[Standard quantified statement]
\[
x_n\to-\infty
\Longleftrightarrow
\forall M\in\mathbb{R}\;\exists N\in\mathbb{N}\;\forall n\ge N
\;\bigl(x_n<M\bigr).
\]
\end{remark*}

\begin{remark*}[Definition predicate reading]
\[
\operatorname{DivergesToNegativeInfinity}(x_n)
\coloneqq
\forall M\in\mathbb{R}\;\exists N\in\mathbb{N}\;\forall n\ge N
\;\bigl(x_n<M\bigr).
\]
\end{remark*}

\begin{remark*}[Negated quantified statement]
\[
x_n\not\to-\infty
\Longleftrightarrow
\exists M\in\mathbb{R}\;\forall N\in\mathbb{N}\;\exists n\ge N
\;\bigl(M\le x_n\bigr).
\]
\end{remark*}

\begin{remark*}[Interpretation]
Divergence to $-\infty$ is escape past every lower threshold. It is the
order-dual of divergence to $+\infty$.
\end{remark*}

\begin{remark*}[Dependencies]\hfill
\begin{itemize}
  \item \hyperref[def:sequence]{Sequence}
\end{itemize}
\end{remark*}

% ---------------------------------------------------------
% Definition --- Oscillatory Sequence
% ---------------------------------------------------------

\begin{definitionbox}{Definition (Oscillatory Sequence)}
\begin{definition}[Oscillatory Sequence]
\label{def:oscillatory-sequence}
Let $(x_n)$ be a real sequence. The sequence $(x_n)$ is
\textbf{oscillatory} if it is divergent, does not diverge to $+\infty$,
and does not diverge to $-\infty$.
\end{definition}
\end{definitionbox}

\begin{remark*}[Standard quantified statement]
\[
\begin{aligned}
&(x_n)\text{ is oscillatory}\\
&\quad\Longleftrightarrow\quad
\Bigl[
  \forall L\in\mathbb{R}\;\exists \varepsilon>0\;\forall N\in\mathbb{N}
  \;\exists n\ge N\;\bigl(|x_n-L|\ge\varepsilon\bigr)
\Bigr]\\
&\qquad\land
\Bigl[
  \exists M\in\mathbb{R}\;\forall N\in\mathbb{N}\;\exists n\ge N
  \;\bigl(x_n\le M\bigr)
\Bigr]\\
&\qquad\land
\Bigl[
  \exists M\in\mathbb{R}\;\forall N\in\mathbb{N}\;\exists n\ge N
  \;\bigl(M\le x_n\bigr)
\Bigr].
\end{aligned}
\]
\end{remark*}

\begin{remark*}[Definition predicate reading]
\[
\begin{aligned}
\operatorname{OscillatorySequence}(x_n)
\coloneqq\;&
\operatorname{DivergentSequence}(x_n)\\
&\land\;\neg\operatorname{DivergesToPositiveInfinity}(x_n)\\
&\land\;\neg\operatorname{DivergesToNegativeInfinity}(x_n).
\end{aligned}
\]
\end{remark*}

\begin{remark*}[Interpretation]
Oscillation means failure to settle without one-directional escape. The
sequence may remain bounded, as $(-1)^n$ does, or it may be unbounded while
still repeatedly returning in incompatible directions.
\end{remark*}

\begin{remark*}[Dependencies]\hfill
\begin{itemize}
  \item \hyperref[def:divergent-sequence]{Divergent Sequence}
  \item \hyperref[def:diverges-to-positive-infinity]{Divergence to Positive Infinity}
  \item \hyperref[def:diverges-to-negative-infinity]{Divergence to Negative Infinity}
\end{itemize}
\end{remark*}

% ---------------------------------------------------------
% Theorem --- Divergence to Infinity Implies Real Divergence
% ---------------------------------------------------------

\begin{theorembox}{Theorem (Divergence to Infinity Implies Real Divergence)}
\begin{theorem}[Divergence to Infinity Implies Real Divergence]
\label{thm:divergence-to-infinity-implies-real-divergence}
Let $(x_n)$ be a real sequence. If $x_n\to+\infty$ or $x_n\to-\infty$,
then $(x_n)$ is divergent.

\smallskip
\noindent
\hyperref[prf:divergence-to-infinity-implies-real-divergence]{\textit{Go to proof.}}
\end{theorem}
\end{theorembox}

\begin{remark*}[Standard quantified statement]
\[
\begin{aligned}
&\Bigl[
  \forall M\in\mathbb{R}\;\exists N\in\mathbb{N}\;\forall n\ge N
  \;\bigl(M<x_n\bigr)
\Bigr]\\
&\quad\lor
\Bigl[
  \forall M\in\mathbb{R}\;\exists N\in\mathbb{N}\;\forall n\ge N
  \;\bigl(x_n<M\bigr)
\Bigr]\\
&\qquad\Longrightarrow
\forall L\in\mathbb{R}\;\exists \varepsilon>0\;\forall N\in\mathbb{N}
\;\exists n\ge N\;\bigl(|x_n-L|\ge\varepsilon\bigr).
\end{aligned}
\]
\end{remark*}

\begin{remark*}[Definition predicate reading]
\[
\operatorname{DivergesToPositiveInfinity}(x_n)
\lor
\operatorname{DivergesToNegativeInfinity}(x_n)
\Longrightarrow
\operatorname{DivergentSequence}(x_n).
\]
\end{remark*}

\begin{remark*}[Interpretation]
Infinity divergence is a stronger, directional failure of real convergence.
A sequence escaping past every real threshold cannot approach a finite real
limit.
\end{remark*}

\begin{remark*}[Dependencies]\hfill
\begin{itemize}
  \item \hyperref[def:divergent-sequence]{Divergent Sequence}
  \item \hyperref[def:diverges-to-positive-infinity]{Divergence to Positive Infinity}
  \item \hyperref[def:diverges-to-negative-infinity]{Divergence to Negative Infinity}
\end{itemize}
\end{remark*}

% ---------------------------------------------------------
% Theorem --- Two Subsequential Limits Force Divergence
% ---------------------------------------------------------

\begin{theorembox}{Theorem (Two Subsequential Limits Force Divergence)}
\begin{theorem}[Two Subsequential Limits Force Divergence]
\label{thm:two-subsequential-limits-force-divergence}
Let $(x_n)$ be a real sequence. If $(x_n)$ has two subsequential limits
$L,K\in\mathbb{R}$ with $L\ne K$, then $(x_n)$ is divergent.

\smallskip
\noindent
\hyperref[prf:two-subsequential-limits-force-divergence]{\textit{Go to proof.}}
\end{theorem}
\end{theorembox}

\begin{remark*}[Standard quantified statement]
\[
\begin{aligned}
&\exists L,K\in\mathbb{R}\;
\Bigl(
  L\ne K
  \land
  L\text{ is a subsequential limit of }(x_n)
  \land
  K\text{ is a subsequential limit of }(x_n)
\Bigr)\\
&\quad\Longrightarrow
\forall A\in\mathbb{R}\;\exists \varepsilon>0\;\forall N\in\mathbb{N}
\;\exists n\ge N\;\bigl(|x_n-A|\ge\varepsilon\bigr).
\end{aligned}
\]
\end{remark*}

\begin{remark*}[Definition predicate reading]
\[
\begin{aligned}
&\exists L,K\in\mathbb{R}\;
\bigl(
  L\ne K
  \land \operatorname{SubsequentialLimit}(x_n,L)
  \land \operatorname{SubsequentialLimit}(x_n,K)
\bigr)\\
&\quad\Longrightarrow
\operatorname{DivergentSequence}(x_n).
\end{aligned}
\]
\end{remark*}

\begin{remark*}[Interpretation]
Convergence forces all subsequences to inherit the same limit. Two different
subsequential limits therefore witness an irreparable conflict in the tail
behavior of the original sequence.
\end{remark*}

\begin{remark*}[Dependencies]\hfill
\begin{itemize}
  \item \hyperref[def:subsequential-limit]{Subsequential Limit}
  \item \hyperref[thm:subsequences-preserve-limits]{Subsequences Preserve Limits}
  \item \hyperref[thm:uniqueness-of-limits]{Uniqueness of Limits}
  \item \hyperref[thm:divergence-by-two-subsequential-limits]{Divergence by Two Subsequential Limits}
\end{itemize}
\end{remark*}

% ---------------------------------------------------------
% Theorem --- Unbounded Above Sequence Has a Positive-Infinity Subsequence
% ---------------------------------------------------------

\begin{theorembox}{Theorem (Unbounded Above Sequence Has a Positive-Infinity Subsequence)}
\begin{theorem}[Unbounded Above Sequence Has a Positive-Infinity Subsequence]
\label{thm:unbounded-above-has-positive-infinity-subsequence}
Let $(x_n)$ be a real sequence. If $(x_n)$ is not bounded above, then
there exists a subsequence $(x_{\sigma(k)})$ such that
$x_{\sigma(k)}\to+\infty$.

\smallskip
\noindent
\hyperref[prf:unbounded-above-has-positive-infinity-subsequence]{\textit{Go to proof.}}
\end{theorem}
\end{theorembox}

\begin{remark*}[Standard quantified statement]
\[
\begin{aligned}
&\Bigl[
  \forall M\in\mathbb{R}\;\exists n\in\mathbb{N}\;\bigl(M<x_n\bigr)
\Bigr]\\
&\quad\Longrightarrow
\exists \sigma:\mathbb{N}\to\mathbb{N}\;
\Bigl(
  \forall k,\ell\in\mathbb{N}\;(k<\ell\Longrightarrow \sigma(k)<\sigma(\ell))
  \land
  \forall M\in\mathbb{R}\;\exists K\in\mathbb{N}\;\forall k\ge K
  \;\bigl(M<x_{\sigma(k)}\bigr)
\Bigr).
\end{aligned}
\]
\end{remark*}

\begin{remark*}[Definition predicate reading]
\[
\neg\operatorname{BoundedAboveSequence}(x_n)
\Longrightarrow
\exists \sigma\;
\bigl(
  \operatorname{StrictlyIncreasing}(\sigma)
  \land
  \operatorname{DivergesToPositiveInfinity}(x_{\sigma(k)})
\bigr).
\]
\end{remark*}

\begin{remark*}[Interpretation]
Unboundedness above may be sparse. This theorem says the high values can
still be selected in increasing index order to form a subsequence that
escapes to $+\infty$.
\end{remark*}

\begin{remark*}[Dependencies]\hfill
\begin{itemize}
  \item \hyperref[def:sequence-bounded-above]{Sequence Bounded Above}
  \item \hyperref[def:subsequence]{Subsequence}
  \item \hyperref[def:diverges-to-positive-infinity]{Divergence to Positive Infinity}
\end{itemize}
\end{remark*}

% ---------------------------------------------------------
% Theorem --- Unbounded Below Sequence Has a Negative-Infinity Subsequence
% ---------------------------------------------------------

\begin{theorembox}{Theorem (Unbounded Below Sequence Has a Negative-Infinity Subsequence)}
\begin{theorem}[Unbounded Below Sequence Has a Negative-Infinity Subsequence]
\label{thm:unbounded-below-has-negative-infinity-subsequence}
Let $(x_n)$ be a real sequence. If $(x_n)$ is not bounded below, then
there exists a subsequence $(x_{\sigma(k)})$ such that
$x_{\sigma(k)}\to-\infty$.

\smallskip
\noindent
\hyperref[prf:unbounded-below-has-negative-infinity-subsequence]{\textit{Go to proof.}}
\end{theorem}
\end{theorembox}

\begin{remark*}[Standard quantified statement]
\[
\begin{aligned}
&\Bigl[
  \forall M\in\mathbb{R}\;\exists n\in\mathbb{N}\;\bigl(x_n<M\bigr)
\Bigr]\\
&\quad\Longrightarrow
\exists \sigma:\mathbb{N}\to\mathbb{N}\;
\Bigl(
  \forall k,\ell\in\mathbb{N}\;(k<\ell\Longrightarrow \sigma(k)<\sigma(\ell))
  \land
  \forall M\in\mathbb{R}\;\exists K\in\mathbb{N}\;\forall k\ge K
  \;\bigl(x_{\sigma(k)}<M\bigr)
\Bigr).
\end{aligned}
\]
\end{remark*}

\begin{remark*}[Definition predicate reading]
\[
\neg\operatorname{BoundedBelowSequence}(x_n)
\Longrightarrow
\exists \sigma\;
\bigl(
  \operatorname{StrictlyIncreasing}(\sigma)
  \land
  \operatorname{DivergesToNegativeInfinity}(x_{\sigma(k)})
\bigr).
\]
\end{remark*}

\begin{remark*}[Interpretation]
Unboundedness below can be converted into a selected tail that escapes past
every lower threshold. The proof is the order-reversed version of the
positive-infinity subsequence construction.
\end{remark*}

\begin{remark*}[Dependencies]\hfill
\begin{itemize}
  \item \hyperref[def:sequence-bounded-below]{Sequence Bounded Below}
  \item \hyperref[def:subsequence]{Subsequence}
  \item \hyperref[def:diverges-to-negative-infinity]{Divergence to Negative Infinity}
\end{itemize}
\end{remark*}

% ---------------------------------------------------------
% Theorem --- Bounded Divergence Produces Two Subsequential Limits
% ---------------------------------------------------------

\begin{theorembox}{Theorem (Bounded Divergence Produces Two Subsequential Limits)}
\begin{theorem}[Bounded Divergence Produces Two Subsequential Limits]
\label{thm:bounded-divergence-produces-two-subsequential-limits}
Let $(x_n)$ be a bounded real sequence. If $(x_n)$ is divergent, then
there exist $L,K\in\mathbb{R}$ with $L\ne K$ such that $L$ and $K$ are
subsequential limits of $(x_n)$.

\smallskip
\noindent
\hyperref[prf:bounded-divergence-produces-two-subsequential-limits]{\textit{Go to proof.}}
\end{theorem}
\end{theorembox}

\begin{remark*}[Standard quantified statement]
\[
\begin{aligned}
&\Bigl[
  \exists M>0\;\forall n\in\mathbb{N}\;\bigl(|x_n|\le M\bigr)
\Bigr]
\land
\Bigl[
  \forall A\in\mathbb{R}\;\exists \varepsilon>0\;\forall N\in\mathbb{N}
  \;\exists n\ge N\;\bigl(|x_n-A|\ge\varepsilon\bigr)
\Bigr]\\
&\quad\Longrightarrow
\exists L,K\in\mathbb{R}\;
\Bigl(
  L\ne K
  \land L\text{ is a subsequential limit of }(x_n)
  \land K\text{ is a subsequential limit of }(x_n)
\Bigr).
\end{aligned}
\]
\end{remark*}

\begin{remark*}[Definition predicate reading]
\[
\operatorname{BoundedSequence}(x_n)
\land
\operatorname{DivergentSequence}(x_n)
\Longrightarrow
\exists L,K\in\mathbb{R}\;
\bigl(
  L\ne K
  \land \operatorname{SubsequentialLimit}(x_n,L)
  \land \operatorname{SubsequentialLimit}(x_n,K)
\bigr).
\]
\end{remark*}

\begin{remark*}[Interpretation]
Bounded divergence cannot be explained by escape to infinity. In the real
line, compactness forces convergent subsequences, and divergence forces at
least two incompatible subsequential limits.
\end{remark*}

\begin{remark*}[Dependencies]\hfill
\begin{itemize}
  \item \hyperref[def:bounded-sequence]{Bounded Sequence}
  \item \hyperref[def:divergent-sequence]{Divergent Sequence}
  \item \hyperref[def:subsequential-limit]{Subsequential Limit}
  \item \hyperref[thm:bolzano-weierstrass-sequences]{Bolzano-Weierstrass Theorem for Sequences}
  \item \hyperref[thm:subsequences-preserve-limits]{Subsequences Preserve Limits}
\end{itemize}
\end{remark*}